\documentclass[a4paper,12pt]{article}
\usepackage{graphicx}
\usepackage{hyperref}
\usepackage{geometry}
\usepackage{pdfpages}
\usepackage[ngerman]{babel}
\geometry{margin=1in}
\setlength{\parindent}{0pt}
\usepackage{listings}

\lstdefinestyle{inline}{
  basicstyle=\fontfamily{pcr}\small\bfseries,
  keywordstyle=\color{cyan!75!black}\bfseries,
  commentstyle=\color{green!75!black},
  stringstyle=\color{orange},
  breaklines=true,
}
% Title Page
\title{\textbf{Anleitung zum Raspi-CyberSec-Lab}}
\author{Jonas Schmitt}
\date{\today}

\begin{document}

\maketitle
\vspace{7cm}

\section{Allgemeine Informationen}
Das RaspberryPI Cybersecurity Lab (RCSL) ist eine Platform zum Testen verschiedener Cyberangriffe, das Projekt ist Teil meiner Bachelorarbeit in Elektro- und Informationstechnik an der Technischen Hochschule Nürnberg.
Es hat verschiedene Funktionen um Umgebungen zu schaffen, welche für das Ausprobieren von Angriffen aus den Bereichen Wifi Netzwerke, Webseiten / -applikationen, MQTT Kommunikation und Bluetooth ausgelegt sind.


\section{Benutzeranleitung}

Zum Einschalten des Geräts, schließen Sie das Gerät mit dem zugehörigen Netzteil an den Strom an.
Das Gerät fährt nun hoch, wobei der Bootvorgang zu sehen ist (das RaspberryPI Logo wird angezeigt).
Nach erfolgreichem Bootvorgang wird das Hauptmenü angezeigt.

\newpage

\subsection{Navigation}
Zur Navigation wird das Drehrad, wie in folgender Abbildung beschrieben, genutzt:

\begin{figure}[h]
    \centering 
    \includegraphics[width=0.6\textwidth]{../figures/Abbildungen/navigation.pdf}
    \label{fig:navigation}
\end{figure}


Die farbige Hinterlegung zeigt an, welcher Menüpunkt aktuell ausgewählt ist. 
Mit dem Drücken des Knopfes wird die Auswahl bestätigt und das Untermenü geöffnet, bzw die Aktion ausgeführt.

Zum verlassen eines Menüs wird die Option "back" gewählt.


\subsection{Menü Übersicht}

Dem Nutzer steht der folgende Menübaum zu Verfügung:

\begin{itemize}
    \item \textbf{wifi} - Optionen zu Wifi Netzwerken
        \begin{itemize}
            \item \textbf{activate} - Aktivieren eines Hotspots
                \subitem Auswahl der Netzwerkart
            \item \textbf{deactivate} - Abschalten des aktiven Hotspots
            \item \textbf{status} - Anzeigen der aktuellen Netzwerkeinstellungen
            \item \textbf{change password} - Setzen eines neuen Passworts
                \subitem Auswahl der zu änderenden Netzwerkart
            \item \textbf{monitoring} - Einschalten und Anzeigen des Netzwerkmonitoring
                \subitem on - Einschalten
                \subitem off - Ausschalten
                \subitem show log - Log anzeigen
                \subitem delete log - Log löschen
        \end{itemize}
    \item \textbf{bluetooth} - Optionen zu Bluetooth
    \item \textbf{webapp} - Optionen zu Webapplikationen
        \begin{itemize}
            \item \textbf{Juice Shop}
                \subitem on - Server Einschalten
                \subitem off - Server Ausschalten
            \item \textbf{MQTT}
                \subitem on - MQTT Konversation starten
                \subitem off - MQTT Konversation beenden
        \end{itemize}
    \item \textbf{development} - Anpassbare Optionen für Entwickler
    \item \textbf{power off} - Ausschalten des Geräts
\end{itemize}

%\newpage

\section{Installation}
Das Projekt wurde mit folgender Hardware umgesetzt:
\begin{itemize}
    \item RaspberryPI 4B mit 8GB RAM
    \item Fenvi AX1800 USB Netzwerkkarte
    \item Waveshare ESP32c6 Microcontroller
    \item Waveshare 4.3 zoll LCD 
    \item Drehgeber mit Druckknopf
\end{itemize}

Als Betriebssystem des RaspberryPI dient RaspberryPI OS in der light Version (nur Kommandozeile).
Zum Kompilieren und Flashen des ESP32 Codes wird das ESP-IDF benötigt, dieses kann in Visual Studio Code als Erweiterung installiert werden.\\

 Zuerst muss das Github \href{https://github.com/Der-Erzfeind/Raspi-CyberSec-Lab-Project.git}{Repository} des Projekts auf den RaspberryPI geklont werden, dies kann mit folgendem Kommando erfolgen:
 \begin{verbatim}
    git clone https://github.com/Der-Erzfeind/Raspi-CyberSec-Lab-Project.git
 \end{verbatim}
Es wird empfohlen das Repository in das standardmäßige Nutzerverzeichnis zu klonen.
Folgende Aktionen müssen im Anschluss ausgeführt werden:

\begin{itemize}
    \item Kompilieren von main.cpp, encoder.cpp und mqtt.cpp
    \item Alle Skripte ausführbar machen
    \item Installation und Konfiguration von Mosquitto
    \item Installation des JuiceShops (und node.js)
    \item Einrichten des systemd services zum Ausführen von start.sh beim Start des Geräts
    \item Kompilieren und Flashen des ESP32 Codes
\end{itemize}
%Nach dem Herunterladen der Projektdaten muss das Installationsskript setup.sh ausgeführt werden.

\section{Entwicklung}
Das RCSL öffnet nach dem Bootvorgang einen Hotspot, der für den Aufbau einer SSH Verbindung, wie folgend, genutzt werden kann:
\begin{enumerate}
    \item Mit dem Hotspot verbinden, \\SSID: RPI \\Passwort: CyberSec
    \item SSH Verbindung aufbauen mit \lstinline[style=inline]|ssh pi@10.40.0.1| \\Passwort: admin
\end{enumerate}
Im Nutzerverzeichnich befinden sich die Ordner mit den Daten zum Projekt und dem Juice Shop.
In der Datei \textit{~/Raspi-CyberSec-Lab-Project/Skripte/devMenu.sh} können die Entwicklungsfunktionen, welche im Menü "development" verfügbar sind, angepasst werden.

\end{document}